% !TEX program = xelatex
\documentclass[tikz,border=8pt]{standalone}
\usepackage{iftex}
\ifPDFTeX
  \errmessage{This figure requires XeLaTeX. In Overleaf, switch Menu -> Compiler -> XeLaTeX.}
\fi
\usepackage{fontspec}
\usepackage{xeCJK}
\usepackage{amsmath}
\usetikzlibrary{arrows.meta,calc}

\setmainfont{TeX Gyre Termes}
\IfFontExistsTF{Noto Serif CJK SC}
  {\setCJKmainfont{Noto Serif CJK SC}}
  {
    \IfFontExistsTF{Source Han Serif SC}
      {\setCJKmainfont{Source Han Serif SC}}
      {
        \IfFontExistsTF{FandolSong-Regular}
          {\setCJKmainfont{FandolSong-Regular}}
          {\setCJKmainfont{SimSun}}
      }
  }

\definecolor{TitleBlue}{RGB}{24,35,220}
\definecolor{LayerTeal}{RGB}{51,158,164}
\definecolor{LayerOrange}{RGB}{255,145,47}
\definecolor{LayerGreen}{RGB}{0,168,0}
\definecolor{LayerBlue}{RGB}{25,145,232}
\definecolor{LayerOrangeDeep}{RGB}{255,120,24}

\definecolor{BoxTealBorder}{RGB}{60,165,171}
\definecolor{BoxOrangeBorder}{RGB}{255,145,47}
\definecolor{BoxGreenBorder}{RGB}{0,168,0}
\definecolor{BoxBlueBorder}{RGB}{25,145,232}
\definecolor{BoxOrangeDeepBorder}{RGB}{255,120,24}

\tikzset{
  titlebox/.style={
    fill=TitleBlue,
    rounded corners=1pt,
    minimum height=1.30cm,
    text=white,
    font=\bfseries\fontsize{12.2pt}{14.2pt}\selectfont,
    align=center
  },
  side label/.style={
    minimum width=1.30cm,
    text=white,
    font=\bfseries\fontsize{10.4pt}{11.9pt}\selectfont,
    align=center
  },
  tophead/.style={
    minimum width=3.95cm,
    minimum height=0.98cm,
    rounded corners=2pt,
    text=white,
    font=\bfseries\fontsize{10.0pt}{11.8pt}\selectfont,
    align=center,
    text width=3.65cm
  },
  appbox/.style={
    minimum width=3.95cm,
    minimum height=1.78cm,
    rounded corners=2pt,
    fill=white,
    font=\bfseries\fontsize{9.5pt}{11.4pt}\selectfont,
    align=center,
    text width=3.58cm
  },
  controlbox/.style={
    minimum width=4.05cm,
    minimum height=1.28cm,
    rounded corners=2pt,
    fill=white,
    font=\bfseries\fontsize{9.1pt}{11.0pt}\selectfont,
    align=center,
    text width=3.68cm
  },
  featurecard/.style={
    minimum width=4.05cm,
    minimum height=1.34cm,
    rounded corners=2pt,
    fill=white,
    font=\bfseries\fontsize{9.15pt}{11.0pt}\selectfont,
    align=center,
    text width=3.68cm
  },
  protocolcard/.style={
    minimum width=4.10cm,
    minimum height=1.36cm,
    rounded corners=2pt,
    fill=white,
    font=\bfseries\fontsize{9.15pt}{11.0pt}\selectfont,
    align=center,
    text width=3.72cm
  },
  primitivebox/.style={
    minimum width=3.00cm,
    minimum height=1.42cm,
    rounded corners=2pt,
    fill=white,
    font=\bfseries\fontsize{9.35pt}{11.2pt}\selectfont,
    align=center,
    text width=2.62cm
  },
  group title/.style={
    font=\bfseries\fontsize{11.8pt}{13.4pt}\selectfont,
    align=center,
    fill=white,
    inner sep=2.2pt
  },
  group dashed teal/.style={draw=BoxTealBorder, dashed, line width=0.9pt, rounded corners=3pt},
  group dashed orange/.style={draw=BoxOrangeBorder, dashed, line width=0.9pt, rounded corners=3pt},
  group dashed green/.style={draw=BoxGreenBorder, dashed, line width=0.9pt, rounded corners=3pt},
  group dashed blue/.style={draw=BoxBlueBorder, dashed, line width=0.9pt, rounded corners=3pt},
  group dashed orange deep/.style={draw=BoxOrangeDeepBorder, dashed, line width=0.9pt, rounded corners=3pt},
  flow/.style={-Latex, line width=1.0pt, draw=BoxOrangeDeepBorder}
}

\begin{document}
\begin{tikzpicture}[x=1cm,y=1cm]

% ---- Global guides ----
\def\L{1.90}
\def\R{16.50}
\def\C{9.20}

% ---- Title ----
\node[titlebox, minimum width=16.4cm] at (\C, 22.15)
{密捷：面向动态词元剪枝的双向隐私安全推理系统总体架构};

% ---- Left side vertical labels ----
\node[side label, fill=LayerTeal, minimum height=3.72cm]       at (0.95, 18.10) {应\\用\\与\\参\\与\\方};
\node[side label, fill=LayerOrange, minimum height=4.15cm]     at (0.95, 13.55) {控\\制\\面\\闭\\环};
\node[side label, fill=LayerGreen, minimum height=6.15cm]      at (0.95, 8.45)  {模\\型\\与\\推\\理\\语\\义};
\node[side label, fill=LayerBlue, minimum height=4.30cm]       at (0.95, 2.95)  {安\\全\\计\\算\\协\\议};
\node[side label, fill=LayerOrangeDeep, minimum height=3.10cm] at (0.95, -0.40) {密\\码\\学\\与\\审\\计\\基\\础};

% ---- Group 1: 应用与参与方 ----
\draw[group dashed teal] (\L,19.95) rectangle (\R,16.25);

\node[tophead, fill=LayerTeal] at (4.10, 18.98) {医院终端（Client）};
\node[tophead, fill=LayerTeal] at (8.75, 18.98) {业务网关（Server）};
\node[tophead, fill=LayerTeal, minimum width=4.55cm, text width=4.18cm] at (13.50, 18.98) {双节点密态执行域（SPU 2PC）};

\node[appbox, draw=BoxTealBorder, line width=0.9pt] at (4.10, 17.42)
{浏览器上传与预览\\Worker 内完成预处理、DQA 与分片};
\node[appbox, draw=BoxTealBorder, line width=0.9pt] at (8.75, 17.42)
{raw multipart 预检\\结构化元数据校验与重放守卫};
\node[appbox, draw=BoxTealBorder, line width=0.9pt, minimum width=4.55cm, text width=4.18cm] at (13.50, 17.42)
{PredictorLG in-SPU\\两方联合完成密态前向\\仅公开最终分类结果};

% ---- Group 2: 控制面闭环 ----
\draw[group dashed orange] (\L,16.10) rectangle (\R,11.95);

\node[controlbox, draw=BoxOrangeBorder, line width=0.9pt] at (4.10, 14.95)
{浏览器 Worker\\图像头部嗅探};
\node[controlbox, draw=BoxOrangeBorder, line width=0.9pt] at (9.20, 14.95)
{浏览器 Worker\\中心裁剪与归一化};
\node[controlbox, draw=BoxOrangeBorder, line width=0.9pt] at (14.30, 14.95)
{浏览器 Worker\\本地 DQA 与审计哈希链};

\node[controlbox, draw=BoxOrangeBorder, line width=0.9pt] at (4.10, 13.08)
{little-endian\\share 序列化};
\node[controlbox, draw=BoxOrangeBorder, line width=0.9pt] at (9.20, 13.08)
{服务端权威快检\\JSON 字节门与 share 哈希校验};
\node[controlbox, draw=BoxOrangeBorder, line width=0.9pt] at (14.30, 13.08)
{NaN/Inf/次正规数阻断\\重放守卫、限频与审计落盘};

% ---- Group 3: 模型与推理语义 ----
\draw[group dashed green] (\L,11.35) rectangle (\R,5.20);
\node[group title, text=BoxGreenBorder] at (\C, 10.78)
{DynamicViT 安全化重写与 MPC-Friendly Transformer 组件};

\node[featurecard, draw=BoxGreenBorder, line width=0.9pt] at (4.10, 9.48)
{pruning boundary 重写\\masking $\rightarrow F_{mux}$};
\node[featurecard, draw=BoxGreenBorder, line width=0.9pt] at (9.20, 9.48)
{threshold compare\\$\rightarrow F_{less}$};
\node[featurecard, draw=BoxGreenBorder, line width=0.9pt] at (14.30, 9.48)
{安全 Top-K\\与 tie 处理};

\node[featurecard, draw=BoxGreenBorder, line width=0.9pt] at (4.10, 7.75)
{Encoded-Key\\Bitonic Sort};
\node[featurecard, draw=BoxGreenBorder, line width=0.9pt] at (9.20, 7.75)
{安全友好算子族\\exact LayerNorm / uniform attention / fixed\_square};
\node[featurecard, draw=BoxGreenBorder, line width=0.9pt] at (14.30, 7.75)
{正式交付主线\\医疗主线\\金融仅作压力验证};

\node[featurecard, draw=BoxGreenBorder, line width=0.9pt] at (4.10, 6.02)
{按索引稳定构造\\keep-mask};
\node[featurecard, draw=BoxGreenBorder, line width=0.9pt] at (9.20, 6.02)
{动态 secure pruning\\保留 full privacy};
\node[featurecard, draw=BoxGreenBorder, line width=0.9pt] at (14.30, 6.02)
{边界压力验证\\dynamic routing 稳定性};

% ---- Group 4: 安全计算协议 ----
\draw[group dashed blue] (\L,4.90) rectangle (\R,0.95);
\node[group title, text=BoxBlueBorder] at (\C, 4.35)
{两方半诚实安全计算协议};

\node[protocolcard, draw=BoxBlueBorder, line width=0.9pt] at (4.10, 2.85)
{线性协议\\矩阵乘与线性层};
\node[protocolcard, draw=BoxBlueBorder, line width=0.9pt] at (9.20, 2.85)
{归一化相关计算\\安全比较 $F_{less}$};
\node[protocolcard, draw=BoxBlueBorder, line width=0.9pt] at (14.30, 2.85)
{安全选择 $F_{mux}$\\排序与决策链};

% ---- Group 5: 密码学与审计基础 ----
\draw[group dashed orange deep] (\L,0.35) rectangle (\R,-2.75);
\node[group title, text=BoxOrangeDeepBorder] at (\C, -0.18)
{密码学与审计基础原语};

\node[primitivebox, draw=BoxOrangeDeepBorder, line width=0.9pt] at (4.05, -1.48) {加法秘密分享};
\node[primitivebox, draw=BoxOrangeDeepBorder, line width=0.9pt] at (7.50, -1.48) {Beaver 三元组};
\node[primitivebox, draw=BoxOrangeDeepBorder, line width=0.9pt] at (10.95, -1.48) {布尔/算术\\比较与选择};
\node[primitivebox, draw=BoxOrangeDeepBorder, line width=0.9pt] at (14.40, -1.48) {SHA-256 / BLAKE2s\\审计摘要};

% ---- Caption ----
\node[font=\fontsize{12pt}{14pt}\selectfont] at (\C, -3.18)
{图 1：密捷项目当前框架总体架构图};

\end{tikzpicture}
\end{document}
